\section{Numerical results and resonance analysis}
\label{sec:results}

The formalism developed above expresses the differential conductance in
terms of the microscopic Andreev-reflection amplitude.  In the subgap
regime, \(|E|<\Delta_0\), no propagating quasiparticle mode exists in the
superconducting lead, and the conductance reduces to
\begin{equation}
	G_A^{(g)}(E,N)
	=
	\frac{1}{2\pi}
	\sum_{\sigma=\pm1}
	\int_{-\pi}^{\pi}
	dK\,
	\chi_{\rm in}(E,K,\sigma)\,
	2
	\left|
	r_{he}^{(g)}(E,K,\sigma,N)
	\right|^2 .
	\label{eq:results_subgap_conductance}
\end{equation}
The purpose of this section is to evaluate Eq.~\eqref{eq:results_subgap_conductance}
for the two contact geometries and to relate the resulting conductance
ridges to the round-trip matrix
\begin{equation}
	\mathcal M_{\rm rt}^{(g)}
	=
	R_F P_-(N) R_S^{(g)} P_+(N).
	\label{eq:results_round_trip_matrix}
\end{equation}

\subsection{Numerical parameters and consistency checks}
\label{subsec:numerical_parameters}

Unless otherwise stated, all energies are measured in units of the
intralayer hopping \(t\).  The interlayer hopping is denoted by \(t_\perp\),
the superconducting gap by \(\Delta_0\), and the exchange field in the
ferromagnetic lead by \(h\).  The central bilayer contains \(N\) longitudinal
cells, and the transverse momentum \(K\in[-\pi,\pi]\) is integrated by
numerical quadrature.

For each point \((E,K,\sigma,N)\), we compute the reflection and transmission
probabilities
\begin{equation}
	R_N=
	|r_{ee}|^2,
	\qquad
	R_A=
	|r_{he}|^2,
	\qquad
	T_{\rm QP}
	=
	\sum_{\nu\in{\rm prop.}}
	|t_\nu|^2 .
\end{equation}
The numerical solution is accepted only when the flux-conservation residual
\begin{equation}
	\delta_{\rm flux}
	=
	\left|
	1-R_N-R_A-T_{\rm QP}
	\right|
	\label{eq:flux_residual}
\end{equation}
remains close to machine precision over the open-channel region
\(\chi_{\rm in}=1\).

% Suggested text once parameters are fixed:
% In the calculations below we use
% \(t_\perp/t=\cdots\), \(\Delta_0/t=\cdots\),
% \(\mu_F/t=\cdots\), \(\mu_C/t=\cdots\), \(\mu_S/t=\cdots\),
% and \(h/t=\cdots\).  The transverse integral is evaluated using
% \(N_K=\cdots\) points.

\subsection{Andreev conductance maps}
\label{subsec:andreev_maps}

Figure~\ref{fig:andreev_maps} shows the subgap Andreev conductance
\(G_A^{(g)}(E,N)\) for the two contact geometries.  The maps display a
sequence of sharp ridge-like structures as a function of the excitation
energy \(E\) and the bilayer length \(N\).  These ridges are the most
prominent finite-size effect of the normal bilayer region.

\begin{figure*}[t]
	\centering
	\includegraphics[width=\textwidth]{figures/results/andreev_maps.pdf}
	\caption{
	Subgap Andreev conductance maps for the two contact geometries.
	Panel (a) shows \(G_A^{(1\to1)}(E,N)\), while panel (b) shows
	\(G_A^{(1\to2)}(E,N)\).  The color scale represents the dimensionless
	conductance per transverse unit cell, in units of \(e^2/h\).
	The ridge-like structures originate from coherent round trips inside
	the finite bilayer region.
	}
	\label{fig:andreev_maps}
\end{figure*}

The two geometries share the same central bilayer spectrum, but differ in
the right-interface selector matrices and therefore in the superconducting
reflection matrix \(R_S^{(g)}\).  Consequently, the ridge pattern is not
determined by the bilayer dispersion alone.  It also depends on how each
central mode couples to the superconducting lead.

\subsection{Fixed-\texorpdfstring{\(K\)}{K} channel resonances}
\label{subsec:fixed_k_resonances}

The integrated conductance \(G_A(E,N)\) is obtained only after summing over
spin and integrating over transverse momentum.  The microscopic resonance
condition, however, is naturally formulated at fixed \((K,\sigma)\).  We
therefore first examine the channel-resolved conductance
\begin{equation}
	g_A^{(g)}(E,K,\sigma,N)
	=
	\chi_{\rm in}(E,K,\sigma)
	2
	\left|
	r_{he}^{(g)}(E,K,\sigma,N)
	\right|^2 .
	\label{eq:fixed_k_conductance}
\end{equation}

\begin{figure}[t]
	\centering
	\includegraphics[width=\columnwidth]{figures/results/fixed_k_resonances.pdf}
	\caption{
		Representative fixed-\(K\) Andreev conductance resonances.  Unlike the
		integrated map, this quantity probes a single transverse channel and
		therefore exposes the microscopic resonances before transverse averaging.
	}
	\label{fig:fixed_k_resonances}
\end{figure}

At fixed \(K\), sharp maxima of Eq.~\eqref{eq:fixed_k_conductance} can be
traced to the eigenvalues of the round-trip matrix
\(\mathcal M_{\rm rt}^{(g)}(E,K,\sigma,N)\).  This fixed-channel analysis is
the first step toward understanding which microscopic resonances survive in
the integrated conductance.

\subsection{Round-trip eigenvalues and bright resonances}
\label{subsec:round_trip_eigenvalues}

Let
\begin{equation}
	\mathcal M_{\rm rt}^{(g)} v_j^R
	=
	\mu_j v_j^R,
	\qquad
	(v_j^L)^\dagger\mathcal M_{\rm rt}^{(g)}
	=
	\mu_j (v_j^L)^\dagger,
\end{equation}
with biorthogonal normalization
\begin{equation}
	(v_j^L)^\dagger v_k^R=\delta_{jk}.
\end{equation}
For a diagonalizable round-trip matrix, the Andreev amplitude can be written
as
\begin{equation}
	r_{he}^{(g)}
	=
	\sum_j
	\frac{C_j^{(g)}}{1-\mu_j^{(g)}} ,
	\label{eq:rhe_modal_sum_results}
\end{equation}
where
\begin{equation}
	\begin{split}
		C_j^{(g)}
		=&
		\left[
		e_h^\dagger
		T_{F\leftarrow C}
		P_-
		R_S^{(g)}
		P_+
		v_j^R
		\right]
		\\
		&\times
		\left[
			(v_j^L)^\dagger t_{\rm in}
			\right] .
	\end{split}
	\label{eq:residue_definition_results}
\end{equation}
Writing
\begin{equation}
	\mu_j^{(g)}
	=
	\rho_j^{(g)} e^{i\Theta_j^{(g)}},
\end{equation}
a single isolated resonance occurs when
\begin{equation}
	\rho_j^{(g)}\simeq 1,
	\qquad
	\Theta_j^{(g)}\simeq 2\pi\ell,
	\qquad
	\ell\in\mathbb Z .
	\label{eq:phase_condition_results}
\end{equation}
The modulus \(\rho_j\) measures how much of the mode survives one complete
round trip, while \(\Theta_j\) is the accumulated round-trip phase.  The
residue \(C_j\) determines whether the pole is bright or dark in the
physical conductance: a near-zero denominator in Eq.~\eqref{eq:rhe_modal_sum_results}
produces a visible ridge only when the incident electron couples appreciably
to the corresponding round-trip eigenmode and that eigenmode has appreciable
hole readout at the ferromagnetic lead.

A useful real-axis diagnostic of resonance proximity is the smallest singular
value
\begin{equation}
	s_{\min}^{(g)}
	=
	s_{\min}
	\left[
		I_4-\mathcal M_{\rm rt}^{(g)}
		\right].
	\label{eq:smin_results}
\end{equation}
Unlike the determinant, this quantity measures near-singularity directly and
is therefore more robust when several round-trip eigenvalues contribute
simultaneously.

\begin{figure}[t]
	\centering
	\includegraphics[width=\columnwidth]{figures/results/round_trip_diagnostics.pdf}
	\caption{
		Round-trip resonance diagnostics for a representative geometry and spin
		sector.  The plotted quantities may include
		\(s_{\min}(I_4-\mathcal M_{\rm rt})\), \(\rho_j\), \(\Theta_j\), and
		the residue weight \(|C_j|/|1-\mu_j|\).  Visible conductance ridges are
		associated with round-trip eigenvalues close to unity and with non-small
		residues.
	}
	\label{fig:round_trip_diagnostics}
\end{figure}

\subsection{From fixed-\texorpdfstring{\(K\)}{K} resonances to integrated ridges}
\label{subsec:ridge_overlays}

The condition in Eq.~\eqref{eq:phase_condition_results} is a fixed-\(K\)
criterion.  A ridge in the integrated conductance appears only if the
corresponding resonance branch remains sufficiently aligned as \(K\) varies
over the open-channel window.  Let
\begin{equation}
	E_{\ell j}^{(g)}(K,N;\sigma)
\end{equation}
be a solution of
\begin{equation}
	\Theta_j^{(g)}(E,K,\sigma,N)=2\pi\ell .
\end{equation}
A simple measure of transverse dispersion is
\begin{equation}
	\begin{split}
		\delta E_{\ell j}^{(g)}(N;\sigma)
		=&
		\max_{K\in\mathcal W_{\rm in}}
		E_{\ell j}^{(g)}(K,N;\sigma)
		\\
		&-
		\min_{K\in\mathcal W_{\rm in}}
		E_{\ell j}^{(g)}(K,N;\sigma),
	\end{split}
	\label{eq:k_flatness_results}
\end{equation}
where
\begin{equation}
	\mathcal W_{\rm in}(E,\sigma)
	=
	\left\{
	K\in[-\pi,\pi]:
	\chi_{\rm in}(E,K,\sigma)=1
	\right\}.
\end{equation}
Branches with
\begin{equation}
	\delta E_{\ell j}^{(g)}(N;\sigma)
	\lesssim
	\Gamma_{\ell j}^{(g)}(N)
\end{equation}
are expected to survive transverse integration, where
\(\Gamma_{\ell j}^{(g)}(N)\) is the numerical linewidth of the associated
fixed-\(K\) resonance.  Branches that disperse strongly with \(K\) are
washed out by the transverse integral even if they are sharp in a single
channel.

\begin{figure*}[t]
	\centering
	\includegraphics[width=\textwidth]{figures/results/ridge_overlays.pdf}
	\caption{
	Conductance maps with resonance-branch overlays.  The background shows
	\(G_A^{(g)}(E,N)\), while the curves show the predicted branches
	\(E_{\ell j}^{(g)}(N)\) obtained from the round-trip phase condition.
	Agreement between the curves and the bright ridges identifies the ridges
	as coherent finite-length resonances of the bilayer cavity.
	}
	\label{fig:ridge_overlays}
\end{figure*}

This comparison is the central test of the resonance interpretation.  A
visible ridge must satisfy three conditions simultaneously: the round-trip
matrix must be nearly singular, the relevant modal residue must be large
enough to make the pole bright, and the fixed-\(K\) resonance branch must be
flat enough to survive the transverse integral.

\subsection{Comparison between the two geometries}
\label{subsec:geometry_comparison_results}

The two geometries differ only through the superconducting-side interface
maps, but this local change is sufficient to modify the full round-trip
matrix and therefore the resonance pattern.  In the \(1\to1\) geometry, the
ferromagnetic and superconducting monolayers are connected through the same
graphene layer.  In the \(1\to2\) geometry, a quasiparticle injected from the
ferromagnet must couple through the bilayer region before reaching the layer
connected to the superconductor.  The resulting change in \(R_S^{(g)}\) and
in the modal residues \(C_j^{(g)}\) changes which resonances are bright in
the integrated conductance.

% Suggested discussion once figures are available:
% Compare ridge visibility, spacing, curvature, and peak intensity between
% the two maps.  Identify whether one geometry favors longer-lived modes
% (\rho_j closer to one), stronger Andreev conversion, or flatter branches
% in K.

\subsection{Summary of the numerical results}
\label{subsec:results_summary}

The conductance maps reveal that the finite bilayer does not merely act as a
featureless spacer between the ferromagnet and the superconductor.  Instead,
it behaves as a coherent cavity whose internal modes undergo repeated
reflection between the two atomically distinct contacts.  The round-trip
matrix \(\mathcal M_{\rm rt}^{(g)}\) provides the microscopic description of
this cavity.  Its eigenvalues determine the resonance phase and lifetime,
while the residues determine whether the corresponding poles are visible in
the Andreev conductance.  The ridge overlays therefore provide a direct
connection between the transfer-matrix formalism and the structures observed
in \(G_A(E,N)\).
